\section{Profili utente e Personas}
In questo capitolo verranno descritte le diverse tipologie di utenti che interagiranno con l’applicazione. L’analisi dei profili utente costituirà una base fondamentale per le fasi successive del progetto, in particolare per la valutazione euristica e per le attività di riprogettazione dell’interfaccia.
\\
\\
Come detto prima le tipologie di utente sono 3: Admin, Employee e Cliente

\subsection{Admin}
L'admin è colui che può utilizzare l'applicazione per eseguire qualsiasi tipo di procedura in quanto possiede i permessi più alti. 

\subsubsection{Caratteristiche psicologiche}
\begin{table}[h]
	\centering
	\begin{tabular}{|c|c|}
		\hline
		\textbf{Caratteristica} & \textbf{Tipologia} \\
		\hline
		Stile cognitivo & Analitico, verbale, astratto\\
		\hline
		Attitudine & Neutrale \\
		\hline
		Motivazione & Alta \\
		\hline
	\end{tabular}
	\caption{Caratteristiche psicologiche dell'admin}
\end{table}

\subsubsection{Caratteristiche di conoscenza}
\begin{table}[h]
	\centering
	\begin{tabular}{|c|c|}
		\hline
		\textbf{Caratteristica} & \textbf{Tipologia} \\
		\hline
		Livello di alfabetizzazione & Buon livello \\
		\hline
		Titolo di studio & Laurea \\
		\hline
		Abilità dattilografica & Alta \\
		\hline
		Linguaggio nativo & Italiano \\
		\hline
		Alfabetizzazione informatica & Alta \\
		\hline
	\end{tabular}
	\caption{Caratteristiche di conoscenza dell'admin}
\end{table}

\clearpage
\subsubsection{Caratteristiche di esperienza}
\begin{table}[h]
	\centering
	\begin{tabular}{|c|c|}
		\hline
		\textbf{Caratteristica} & \textbf{Tipologia} \\
		\hline
		Esperienza del compito & Esperto \\
		\hline
		Esperienza sulla applicazione & Medio \\
		\hline
		Esperienza dei sistemi informatici & Media \\
		\hline
		Uso di altri sistemi & Scarso \\
		\hline
	\end{tabular}
	\caption{Caratteristiche di esperienza dell'admin}
\end{table}

\subsubsection{Caratteristiche fisiche}
\begin{table}[h]
	\centering
	\begin{tabular}{|c|c|}
		\hline
		\textbf{Caratteristica} & \textbf{Tipologia} \\
		\hline
		Daltonismo & Sì/No \\
		\hline
		Manualità & Variabile \\
		\hline
		Disabilità & Sì/No\\
		\hline
	\end{tabular}
	\caption{Caratteristiche fisiche dell'admin}
\end{table}

\subsubsection{Caratteristiche del lavoro e dei compiti}
\begin{table}[h]
	\centering
	\begin{tabular}{|c|c|}
		\hline
		\textbf{Caratteristica} & \textbf{Tipologia} \\
		\hline
		Uso sistema & Obbligatorio \\
		\hline
		Frequenza d'uso & Media \\
		\hline
		Addestramento di base & Obbligatorio formale \\
		\hline
		Categorie di lavoro & Manager \\
		\hline
		Uso di altri strumenti & No \\
		\hline
		Frequenza del turnover & Basso \\
		\hline
		Importanza del compito & Alto \\
		\hline
		Struttura del compito & Media \\
		\hline
	\end{tabular}
	\caption{Caratteristiche del lavoro e dei compiti dell'admin}
\end{table}

\clearpage
\subsection{Employee}
L'employee è colui che lavora per la banca allo sportello. Può svolgere svariati compiti poiché possiede un livello medio di permessi.

\subsubsection{Caratteristiche psicologiche}
\begin{table}[h]
	\centering
	\begin{tabular}{|c|c|}
		\hline
		\textbf{Caratteristica} & \textbf{Tipologia} \\
		\hline
		Stile cognitivo & Variabile \\
		\hline
		Attitudine & Neutrale \\
		\hline
		Motivazione & Alta \\
		\hline
	\end{tabular}
	\caption{Caratteristiche psicologiche dell'employee}
\end{table}

\subsubsection{Caratteristiche di conoscenza}
\begin{table}[h]
	\centering
	\begin{tabular}{|c|c|}
		\hline
		\textbf{Caratteristica} & \textbf{Tipologia} \\
		\hline
		Livello di alfabetizzazione & Livello medio/Buon livello \\
		\hline
		Titolo di studio & Superiore \\
		\hline
		Abilità dattilografica & Alta \\
		\hline
		Linguaggio nativo & Italiano \\
		\hline
		Alfabetizzazione informatica & Moderata \\
		\hline
	\end{tabular}
	\caption{Caratteristiche di conoscenza dell'employee}
\end{table}

\subsubsection{Caratteristiche di esperienza}
\begin{table}[h]
	\centering
	\begin{tabular}{|c|c|}
		\hline
		\textbf{Caratteristica} & \textbf{Tipologia} \\
		\hline
		Esperienza del compito & Medio \\
		\hline
		Esperienza sulla applicazione & Medio \\
		\hline
		Esperienza dei sistemi informatici & Media \\
		\hline
		Uso di altri sistemi & Scarso \\
		\hline
	\end{tabular}
	\caption{Caratteristiche di esperienza dell'employee}
\end{table}

\clearpage
\subsubsection{Caratteristiche fisiche}
\begin{table}[h]
	\centering
	\begin{tabular}{|c|c|}
		\hline
		\textbf{Caratteristica} & \textbf{Tipologia} \\
		\hline
		Daltonismo & Sì/No \\
		\hline
		Manualità & Variabile \\
		\hline
		Disabilità & Sì/No \\
		\hline
	\end{tabular}
	\caption{Caratteristiche fisiche dell'employee}
\end{table}

\subsubsection{Caratteristiche del lavoro e dei compiti}
\begin{table}[h]
	\centering
	\begin{tabular}{|c|c|}
		\hline
		\textbf{Caratteristica} & \textbf{Tipologia} \\
		\hline
		Uso sistema & Obbligatorio \\
		\hline
		Frequenza d'uso & Alta \\
		\hline
		Addestramento di base & Obbligatorio formale \\
		\hline
		Categorie di lavoro & Impiegato \\
		\hline
		Uso di altri strumenti & Attrezzature d'ufficio \\
		\hline
		Frequenza del turnover & Basso \\
		\hline
		Importanza del compito & Alto \\
		\hline
		Struttura del compito & Media \\
		\hline
	\end{tabular}
	\caption{Caratteristiche del lavoro e dei compiti dell'employee}
\end{table}

\clearpage
\subsection{Cliente}
Il cliente è colui che vuole usufruire dei servizi della banca da casa. La clientela può essere molto variegata in quanto chiunque potrebbe voler usufruire dei servizi online.

\subsubsection{Caratteristiche psicologiche}
\begin{table}[h]
	\centering
	\begin{tabular}{|c|c|}
		\hline
		\textbf{Caratteristica} & \textbf{Tipologia} \\
		\hline
		Stile cognitivo & Variabile\\
		\hline
		Attitudine & Variabile \\
		\hline
		Motivazione & Alta \\
		\hline
	\end{tabular}
	\caption{Caratteristiche psicologiche del cliente}
\end{table}

\subsubsection{Caratteristiche di conoscenza}
\begin{table}[h]
	\centering
	\begin{tabular}{|c|c|}
		\hline
		\textbf{Caratteristica} & \textbf{Tipologia} \\
		\hline
		Livello di alfabetizzazione & Livello Medio/Buon Livello \\
		\hline
		Titolo di studio & Variabile \\
		\hline
		Abilità dattilografica & Media/Alta \\
		\hline
		Linguaggio nativo & Variabile \\
		\hline
		Alfabetizzazione informatica & Variabile \\
		\hline
	\end{tabular}
	\caption{Caratteristiche di conoscenza del cliente}
\end{table}

\subsubsection{Caratteristiche di esperienza}
\begin{table}[h]
	\centering
	\begin{tabular}{|c|c|}
		\hline
		\textbf{Caratteristica} & \textbf{Tipologia} \\
		\hline
		Esperienza del compito & Nuovo nel campo/Medio \\
		\hline
		Esperienza sulla applicazione & Novizio/Medio \\
		\hline
		Esperienza dei sistemi informatici & Variabile \\
		\hline
		Uso di altri sistemi & Scarso \\
		\hline
	\end{tabular}
	\caption{Caratteristiche di esperienza del cliente}
\end{table}

\clearpage
\subsubsection{Caratteristiche fisiche}
\begin{table}[h]
	\centering
	\begin{tabular}{|c|c|}
		\hline
		\textbf{Caratteristica} & \textbf{Tipologia} \\
		\hline
		Daltonismo & Sì/No \\
		\hline
		Manualità & Variabile \\
		\hline
		Disabilità & Sì/No\\
		\hline
	\end{tabular}
	\caption{Caratteristiche fisiche del cliente}
\end{table}

\subsubsection{Caratteristiche del lavoro e dei compiti}
\begin{table}[h]
	\centering
	\begin{tabular}{|c|c|}
		\hline
		\textbf{Caratteristica} & \textbf{Tipologia} \\
		\hline
		Uso sistema & Facoltativo \\
		\hline
		Frequenza d'uso & Bassa \\
		\hline
		Addestramento di base & Nessuno \\
		\hline
		Categorie di lavoro & Variabile \\
		\hline
		Uso di altri strumenti & No \\
		\hline
		Frequenza del turnover & Variabile \\
		\hline
		Importanza del compito & Alto \\
		\hline
		Struttura del compito & Media \\
		\hline
	\end{tabular}
	\caption{Caratteristiche del lavoro e dei compiti del cliente}
\end{table}

\clearpage
\subsection{Personas}